\subsection*{Purpose of this Document}
This Governance Charter defines how ScienceEcosystem is stewarded, how decisions are made, and how independence, transparency, and long-term integrity are safeguarded. It is designed to function from the earliest founder-led stage through future institutional maturity.

This document establishes constraints on authority, not merely roles.

\section*{1. Nature of ScienceEcosystem}

ScienceEcosystem is a public-interest digital infrastructure created to support open, verifiable, and reproducible science across generations. It is not a commercial product, publisher, or analytics platform.

ScienceEcosystem exists to serve the scientific community as a whole and is governed accordingly.

\section*{2. Core Governance Principles}

All governance within ScienceEcosystem is guided by the following non-negotiable principles:

\begin{itemize}
    \item \textbf{Public Interest First:} Decisions must prioritise scientific integrity and societal benefit over institutional, personal, or financial interests.
    \item \textbf{Non-Commercial Independence:} ScienceEcosystem shall not be sold, acquired, or converted into a for-profit entity.
    \item \textbf{Transparency:} Decision-making processes, governance structures, and financial information shall be publicly documented.
    \item \textbf{Accountability:} Authority must be reviewable, limited, and transferable.
    \item \textbf{Durability:} Governance choices must favour long-term stability over short-term optimisation.
\end{itemize}

\section*{3. Stewardship Model}

\subsection*{3.1 Founder Stewardship (Initial Phase)}

ScienceEcosystem is currently stewarded by its founder, who acts as:
\begin{itemize}
    \item Initial maintainer of the platform
    \item Custodian of its mission and principles
    \item Interim decision-maker for technical and organisational matters
\end{itemize}

Founder stewardship is explicitly temporary. Its purpose is to enable early development while preparing for institutional governance.

\subsection*{3.2 Limits of Founder Authority}

The founder may not unilaterally:
\begin{itemize}
    \item Introduce paywalls or monetise access to content
    \item Sell user data or behavioural analytics
    \item Enter exclusive partnerships that restrict openness
    \item Alter core governance principles without a public process
\end{itemize}

\section*{4. Governance Bodies}

As ScienceEcosystem matures, governance authority will be distributed across the following bodies.

\subsection*{4.1 Board of Trustees}

\textbf{Role:}
\begin{itemize}
    \item Fiduciary oversight
    \item Protection of mission and independence
    \item Approval of budgets and strategic direction
\end{itemize}

\textbf{Composition:}
\begin{itemize}
    \item Independent members with expertise in science, infrastructure, law, and ethics
    \item Fixed terms with staggered rotation
\end{itemize}

\subsection*{4.2 Scientific Council}

\textbf{Role:}
\begin{itemize}
    \item Advise on scientific priorities and disciplinary balance
    \item Review platform alignment with research practices
\end{itemize}

Members are appointed based on disciplinary expertise and community standing.

\subsection*{4.3 Community Council}

\textbf{Role:}
\begin{itemize}
    \item Represent active users and contributors
    \item Advise on usability, accessibility, and community norms
\end{itemize}

Members are elected by the contributor community.

\subsection*{4.4 Ethics and Integrity Board}

\textbf{Role:}
\begin{itemize}
    \item Safeguard neutrality, inclusivity, and integrity
    \item Review conflicts of interest
    \item Oversee ethical complaints and moderation processes
\end{itemize}

\section*{5. Decision-Making Processes}

\subsection*{5.1 Transparency}

All significant decisions must be:
\begin{itemize}
    \item Documented publicly
    \item Accompanied by a rationale
    \item Subject to community feedback where appropriate
\end{itemize}

\subsection*{5.2 Change Management}

Changes to:
\begin{itemize}
    \item Governance principles
    \item Data policies
    \item Funding models
\end{itemize}

require a public consultation period and formal approval by the Board of Trustees once established.

\section*{6. Funding and Conflict of Interest}

ScienceEcosystem may accept funding only if it:
\begin{itemize}
    \item Does not compromise independence
    \item Does not create exclusive control or influence
    \item Is publicly disclosed
\end{itemize}

All trustees, council members, and senior contributors must disclose conflicts of interest annually.

\section*{7. Open Infrastructure and Forkability}

ScienceEcosystem commits to open standards and open-source infrastructure wherever feasible. This ensures that:
\begin{itemize}
    \item The platform remains auditable
    \item No single entity can capture or enclose it
    \item The community retains a credible exit option
\end{itemize}

Forkability is treated as a governance safeguard, not a failure mode.

\section*{8. Accountability and Reporting}

ScienceEcosystem shall publish:
\begin{itemize}
    \item Annual activity and governance reports
    \item Financial summaries
    \item Roadmap updates
\end{itemize}

These reports are part of its public accountability obligations.

\section*{9. Governance Evolution}

This Charter is a living document. Governance structures will evolve as ScienceEcosystem grows, but changes must:
\begin{itemize}
    \item Strengthen independence
    \item Increase accountability
    \item Preserve public trust
\end{itemize}

At no stage may governance evolution reduce transparency or consolidate unchecked authority.

\section*{10. Commitment}

ScienceEcosystem is stewarded with the understanding that its value lies not in ownership, but in trust. Governance exists to ensure that this trust can survive individual involvement and endure across generations.


